% Flavor catalog per-process draft.
% process_id: L002
% family: charged_lepton
% owner_agent_id: PKA-L002
% writer_agent_id: null
% checker_agent_id: null
% opus_signoff_id: null
% Canonical metadata sidecar: L002.yaml
% Status is controlled by the YAML sidecar status_history, not this comment.

\section{L002: \texorpdfstring{\(\mu\to 3e\)}{mu -> 3e} Four-Lepton LFV}

\subsection*{Process}
\textbf{Standard notation:} \(\mathcal{B}(\mu^+\to e^+e^-e^+)\).
This entry covers the neutrinoless charged-lepton-flavor-violating
three-electron muon decay, with emphasis on four-lepton/contact operators and
the off-shell dipole contribution.  It is a charged-lepton process and is
separate from the quark \(\Delta F=2\) scan lane.

\subsection*{PDG-or-equivalent value}
The current PDG 2026 muon listing, datablock S004R4, gives the sidecar
\texttt{L002.yaml:primary\_current\_limit} value
\[
  \mathcal{B}(\mu^+\to e^+e^-e^+) < 1.0\times 10^{-12}\quad (90\%\ {\rm C.L.}).
\]
PDG attributes the limit to the 1988 SINDRUM spectrometer result of Bellgardt
et al.; the same value is recorded in \texttt{L002.yaml:original\_experiment}.
The local PDG snapshot is
\texttt{flavor\_catalog/references/L002/pdg2026\_muon\_listing\_s004r4.txt};
the SINDRUM metadata snapshot is
\texttt{flavor\_catalog/references/L002/sindrum1988\_inspire\_abstract.txt}.
This remains an upper limit, not a measurement.  The sidecar \texttt{prospects}
entries keep Mu3e separate from the present bound: the phase-I technical design
targets a single-event sensitivity of \(2\times10^{-15}\), while the broader
program aims at \(\mathcal{O}(10^{-16})\) after future upgrades.

\subsection*{Relevance to RS / anarchic flavor}
In a warped lepton extension, \(\mu\to 3e\) probes flavor-violating neutral
currents and four-lepton operators generated by KK gauge exchange, \(Z\)-like
mixing, or other lepton-current misalignment.  It is complementary to L001:
\(\mu\to e\gamma\) is dominantly dipole-like in the current repo, while
\(\mu\to 3e\) can be tree-level/contact dominated and therefore has different
dependence on lepton localization and five-dimensional Yukawa assumptions.
The Agashe--Blechman--Petriello RS LFV analysis already treated rare muon decays
as part of the lepton-sector flavor test set.

\subsection*{Post-2008 developments}
Relative to the arXiv:0804.1954 RS-flavor baseline, the experimental limit has
not yet been superseded in PDG; SINDRUM still sets the catalog value.  The main
post-2008 change is the Mu3e program at PSI: the phase-I detector design,
published in 2021, is built for \(10^8\) muon decays per second and a
\(2\times10^{-15}\) single-event sensitivity
(\texttt{L002.yaml:prospects}).  The 2025 status note records a phase-I
sensitivity scale of \(\mathcal{O}(10^{-15})\) and an upgraded scale of
\(\mathcal{O}(10^{-16})\).  On the theory side, the sidecar source
\texttt{crivellin2017\_mu\_e\_eft\_arxiv1702\_03020.txt} covers systematic
EFT/RG analyses of \(\mu\to e\gamma\), \(\mu\to 3e\), and coherent conversion
below \(m_W\), including QED/QCD operator mixing.

\subsection*{Constraint validity / model dependence}
The experimental interpretation is robust: the Standard-Model rate with massive
neutrinos is negligible, quoted by Mu3e status material as
\(\mathcal{O}(10^{-54})\), so an observation would be new physics.  The mapping
to this repo is model-dependent because a four-lepton Wilson basis, dipole
interference convention, and possible RG evolution must be chosen.  This entry
should be classified as lepton-extension-only and contact-operator sensitive,
not as a quark-scan constraint.

\subsection*{Code coverage in this repo}
\textbf{NO.}  The required greps over \texttt{quarkConstraints/}, \texttt{qcd/},
\texttt{flavorConstraints/}, \texttt{neutrinos/}, \texttt{yukawa/},
\texttt{warpConfig/}, \texttt{solvers/}, \texttt{scanParams/}, and
\texttt{tests/} found no \(\mu\to 3e\), Mu3e, or four-lepton-contact
implementation.  The only adjacent charged-lepton LFV path is the dipole
\(\mu\to e\gamma\) checker at \texttt{flavorConstraints/muToEGamma.py}:75 and
its scan call at \texttt{scanParams/scan.py}:524, which do not evaluate
\(\mu\to 3e\).

\subsection*{Implementation difficulty}
\textbf{MEDIUM.}  A minimal implementation needs a new low-energy
\(\mu\to 3e\) observable formula, a four-lepton Wilson/operator convention, and
a matching hook from the RS lepton-sector parameters.  It should not require
lattice matrix elements or hadronic long-distance inputs.  A production EFT
treatment with QED/QCD RG mixing and correlated \(\mu\to e\gamma\),
\(\mu\to 3e\), and \(\mu\)-to-\(e\) conversion fits would raise the integration
difficulty.

\subsection*{Key references}
Process-local source keys before bibliography consolidation:
\texttt{PDG2026\_MuonS004R4}, \texttt{SINDRUM1988\_Mu3e},
\texttt{Mu3eTDR2021}, \texttt{Mu3eStatus2025},
\texttt{CrivellinDavidsonPrunaSigner2017\_MuEFT},
\texttt{AgasheBlechmanPetriello2006\_RSLFV}, and \texttt{CFW2008}.
